\documentclass[10pt,a4paper]{amsart}
\usepackage[margin=19mm]{geometry}
\usepackage{amsmath,amssymb,mathtools}
\usepackage[scr=rsfs,cal=euler]{mathalfa}
\usepackage{xfrac}
\usepackage[unicode,hidelinks,bookmarks=false]{hyperref}
\newcommand{\ie}{i.e.\ }
\newcommand{\cf}{cf.\ }
\newcommand{\resp}{respectively\ }
\newcommand{\Subsection}{\subsection}
\providecommand{\nobreakdash}{\nobreak\hbox{-}}
\setlength{\emergencystretch}{2em}
\multlinegap0pt
\usepackage{amssymb}
\usepackage{xcolor}
\usepackage{mathtools}
\usepackage{enumerate}
\usepackage[shortlabels]{enumitem}
\newtheorem{theorem}{Theorem}
\newtheorem{claim}{Claim}
\usepackage{cleveref}
\crefname{theorem}{Theorem}{Theorems}
\crefname{claim}{Claim}{Claims}
\title{Unit distance graphs with few crossings per edge}
\author{Panna Geh\'er, D\"om\"ot\"or P\'alv\"olgyi, D\'aniel G. Simon, G\'eza T\'oth}
\date{Source: arXiv:2603.19848v1 (CC BY 4.0)}
\begin{document}
\maketitle

\noindent\emph{Source and licence.} Unit distance graphs with few crossings per edge. Version arXiv:2603.19848v1 (CC BY 4.0), distributed by arXiv under the Creative Commons Attribution 4.0 International licence, http://creativecommons.org/licenses/by/4.0/, as recorded in the arXiv licence metadata for this exact version and shown on the abstract page https://arxiv.org/abs/2603.19848v1. Full licence notice: \texttt{licenses/arxiv-2603.19848-CC-BY-4.0.txt}.\par\smallskip

\noindent\textit{Editorial scope.} Counted unit: the article's theorem thm\_construction (source0.tex lines 128--132). The article's complete original proof of it is delivered with it (source0.tex lines 678--771, one \texttt{proof} environment of 94 lines, which the article places away from the statement). No mathematics is written, repaired or re-derived: the statement and the proof are the article's own lines, in the article's own notation, and no retained line is substituted or re-flowed.  No further source-local context is needed: every cross-reference of every retained line resolves inside the two delivered spans.  One declared adaptation: exactly 3 ancillary strings inside the retained spans are omitted (image-inclusion commands and, where the source repeats a label, the repeated label definitions) and no other line differs from the source; the figure environments, with their placement, captions and labels, are kept, and the package records the omitted strings in fidelity receipts bound to the affected bodies.  No retained line contains a citation, so the wrapper carries no bibliography and no bibliography entry.  The retained lines contain the article's own diagram code, which the wrapper typesets with the standard tikz packages; no image file is delivered and the package declares no asset dependency.  Editorial formatting only, in addition: an AMS \texttt{amsart} preamble that restates the article's own declarations and macros used by the retained lines, namely the environments \texttt{theorem}, \texttt{claim} on separate counters, together with the standard packages those lines need.  The retained environments print in this document as Theorem 1, Claim 1 in order of appearance, whereas the article prints the same items with the numbering of its own full document.  The complete native source of the article as distributed in the arXiv e-print for this exact version is bundled unchanged as \texttt{sources/196804/source0.tex}.
\begin{theorem} \label{thm_construction}
For all $n$ large enough, we have
    $$u_2(n) \geq u_0(n)+0.58\sqrt{n}.$$
    
\end{theorem}

\begin{proof}[\bf Proof of \Cref{thm_construction}.]

Suppose first that $n=69k^2-57k+17$ for some integer $k$. We begin by describing our construction in this special case.
The base unit is a 2-planar unit distance graph $H$ with $9$ vertices and $18$ edges, shown in \Cref{n=9}. This graph is the $3\times 3$ Rook's graph (it is also the Paley graph of order 9 and the Hamming graph $H(2,3)$).
This graph has the most edges among all unit distance graphs on $9$ vertices.

    \begin{figure}[ht!]
    \centering
        
        \caption{The basis of the construction. It has $9$ vertices and $18$ edges.}
        \label{n=9}
    \end{figure}   
    
We build our construction using $H$ as unit blocks. First, using four rotated copies of $H$ and four extra edges, we create a regular dodecagon, as shown in \Cref{n=29}, where the additional edges are shown with dashed lines. Clearly, it is a $2$-planar unit distance graph. It contains $29$ vertices and $72$ edges.


       \begin{figure}[ht!]
      \centering
        
        \caption{The regular dodecagon. We rotated the original $9$-point construction (drawn with thick edges) around its lower left vertex to create a dodecagon. We also added four extra edges (drawn with dashed lines). The graph has $29$ vertices and $72$ edges.}
        \label{n=29}
    \end{figure}
    
Next, we use these dodecagons to form a piece of a lattice. We join the dodecagons in a way such that every second side of the dodecagon is touching another one. In other words, we build a hexagonal grid using the dodecagons. We build the grid following a spiral line, starting from the center. 
Let $k$ denote the number of layers, that is, if the grid contains one dodecagon, we have $k=1$, if the grid contains $7$ dodecagons, we have $k=2$, if the grid contains $19$ dodecagons, we have $k=3$, etc. In \Cref{k=2}, we show how the dodecagons are placed in a hexagonal grid via the case $k=2$. Furthermore, between any two neighboring boundary dodecagons, we add an extra edge, colored with gray in the figure.
    
    
       \begin{figure}[ht!]
        \centering
        
        \caption{The $k=2$ case. We join $7$ regular dodecagons in a hexagonal grid pattern. We add one extra edge (drawn as dashed lines) between each pair of neighboring dodecagons on the boundary.}
        \label{k=2}
    \end{figure}    


We denote the resulting graph containing $n$ vertices by $G_n$ and denote the number of edges in $G_n$ by $e$. Clearly, $G_n$ is a 2-planar unit distance graph.
Let us calculate the number of vertices of $G_n$ with $k$ layers. It contains $3k^2-3k+1$ dodecagons, the number of vertices in a dodecagon is $29$, but at any connection two edges coincide. Let $c_1$ denote the number of edge pairs that coincide, and let $h$ denote the number of dodecagons. For $c_1$, we have
    
$$c_1=\frac12 \left[6(3k^2-3k+1)-12k+6 \right]=3h-6k+3.$$
    
Thus, the total number of vertices needed to complete $k$ layers is $$n=29h-2c_1=23h+12k-6=69k^2-57k+17.$$ 
Next, we calculate the number of edges of $G_n$.

\begin{claim} \label{specialcase}
For any $n$ that can be written as $n=69k^2-57k+17$ where $k$ is an integer, the graph $G_n$ contains
$\left\lfloor3n-\sqrt{\frac{192}{23}n-\frac{23088}{529}}\right\rfloor$ edges.
\end{claim}

\begin{proof}
As we have already seen, as $n=69k^2-57k+17$, 
the graph $G_n$ corresponds to the construction with $k$ completed layers.
Each dodecagon contains $72$ edges, at each connection one edge is double-counted, and between any two neighboring dodecagons on the boundary of the grid we added an extra (gray) edge. Let $c_2$ denote the number of neighboring pairs on the boundary. Clearly, we have $c_2=6(k-1)$. Thus,    
$$e =72h-c_1+c_2=72(3k^2-3k+1)-3(3k^2-3k+1)+6k-3+6k-6= 207k^2-195k+60.$$

    As $n=69k^2-57k+17$, we have:
   \begin{align*}
e &=207k^2-195k+60=\left\lfloor207k^2-195k+60+\frac{21}{23}\right\rfloor=\left\lfloor207k^2-171k+51-\sqrt{\left(24k-\frac{228}{23}\right)^2}\right\rfloor
\\
&=\left\lfloor207k^2-171k+51-\sqrt{576k^2-\frac{10944}{23}k+\frac{3264}{23}-\frac{23088}{529}}\right\rfloor=\left\lfloor3n-\sqrt{\frac{192}{23}n-\frac{23088}{529}}\right\rfloor.
    \end{align*}

    Hence, the statement is proved.
\end{proof}


Now we investigate the general case. We can assume that $n$ can be written as $n=69k^2-57k+17+A$ for some integers $k$ and $A<[69(k+1)^2-57(k+1)+17]-[69k^2-57k+17]=138k+12$. 
We extend the construction in the following way. First, we consider the graph $G_m$ for $m=69k^2-57k+17$ and we use the remaining $A$ vertices as follows: we choose a position adjacent to two dodecagons of the previous layer and we begin building a dodecagon. We continue adding vertices to this dodecagon until either the dodecagon is completed, or no vertices remain. Once a dodecagon is fully formed, we proceed to the next position in clockwise order in the $(k+1)$-th layer. We only start building a new dodecagon when the previous one is completed, and repeat this procedure until we place $A$ new vertices.
We still denote the resulting graph containing $n$ vertices by $G_n$. \Cref{thm_construction} follows directly from the following statement.
\begin{theorem}
For all $n \geq 179$, the graph $G_n$ contains
$\left\lfloor3n-\sqrt{\frac{192}{23}n-\frac{23088}{529}}-\frac{13683}{23}\right\rfloor$ edges.
\end{theorem}

\begin{proof}
Let $n=69k^2-57k+17+A$ for some integers $k$ and nonnegative $A<138k+12$ where $n \geq 179$. That is, $G_n$ contains $k \geq 2$ completed layers of the hexagonal grid and an incomplete layer with $A \geq 0$ vertices. By \Cref{specialcase} the $k$ layers contain $207k^2-195k+60$ edges. Consider the remaining $A$ vertices in the $(k+1)$-th layer. When building the $(k+1)$-th layer along the spiral, the first dodecagon requires $25$ extra vertices and the next ones always require either $25$ or $23$ extra vertices, contributing $71$ or $69$ new edges, respectively.
Furthermore, there are only $7$ dodecagons requiring $25$ vertices (the first one and the ones on the corners of the hexagonal grid), which consume at most $7 \cdot 25 = 175$ vertices. For simplicity, we omit these dodecagons. Thus, there are at least $\left\lfloor\frac{A-175}{23}\right\rfloor\geq \frac{A-198}{23}$ fully formed dodecagons in the $(k+1)$-th layer, each contributing $69$ edges.

Thus, for the number of edges of $G_n$, we have
\begin{align}
e&\geq 207k^2-195k+60+\frac{A-198}{23}\cdot69  \nonumber\\
&=207k^2-195k+3A-534 \nonumber\\
&= 3n-24k-585 \label{eq9} \\
&\geq 3n - 24 \left( \frac{57 + \sqrt{276n - 1443}}{138} \right) - 585 \label{eq10}\\
&= 3n - \sqrt{\frac{192}{23}n - \frac{23088}{529}} - \frac{13683}{23}. \nonumber
\end{align}

In \eqref{eq9} we used the fact that $3n=207k^2-171k+51+3A$. In \eqref{eq10}, we used that $n=69k^2-57k+17+A$, $A \geq 0$ and therefore
$k \leq \frac{57 + \sqrt{276n - 1443}}{138}.$

This proves the statement.
\end{proof}

To finish the proof of \Cref{thm_construction}, recall that $u_0(n) = 3n-\left\lfloor\sqrt{12n - 3}\right\rfloor$. Since $\frac{192}{23} \approx 8.34 < 12$, for large enough values of $n$, the number of edges in our construction exceeds $u_0(n)$. The difference is asymptotically $\left(\sqrt{12} - \sqrt{192/23}\right)\sqrt{n} \approx 0.58\sqrt{n}$.
\end{proof}

\end{document}
